\section*{The macroscopic interval likelihood and its approximations}
% Owner: papers/_program/nomenclature.md (letters, lattice, ladder, region names) and
% papers/_program/program.md §1 (body roster). Presents the root question (are the gating
% fluctuations modelled at all?) with least squares at the bottom of the cost ladder, and the
% lattice of likelihoods that hangs from it, with the endpoint ladder as the presentation.
% Do NOT restate: model scope (program.md §7), the novelty claim, the diagnostic definitions
% (03_diagnostics.tex), or run values and the noise parameter (06_methods.tex). The index-level
% kernels for the boundary moments (E_2, E_3, the Kronecker-sum construction) live in Methods and
% the Supplement; this section names them.

A macroscopic current is the summed signal of a population of ion channels, each of
which behaves as an independent continuous-time Markov chain over a small set of
conformational states. We write $N_{\mathrm{ch}}$ for the number of channels and $K$ for
the number of kinetic states of one channel ($K=2$ throughout this paper, carried % src: papers/_program/program.md §7 (model scope: two states)
symbolically where it costs nothing). Each channel is governed by a generator (rate)
matrix $\mathbf{Q}\in\mathbb{R}^{K\times K}$, whose transition matrix over an elapsed time
$t$ is the row-stochastic $\mathbf{P}(t)=e^{\mathbf{Q}t}$ with entries $P_{i\to j}(t)$. A
channel in state $k$ carries a conductance $\gamma_k$, and we collect these in the vector
$\bm{\gamma}\in\mathbb{R}^{K}$.

The likelihood we need is the probability of a recorded sample given the kinetic
parameters and the data seen so far. Its exact form would require the distribution of the
interval-averaged conductance of the whole population, which has no closed expression.
Every practical likelihood replaces that distribution with a Gaussian, and the members of
the family studied here differ in what the Gaussian is conditioned on and in how its
variance is accounted. The rest of this section unpacks that sentence.

\subsection*{The observable is an interval average}

A recorded sample is the current averaged over the acquisition window rather than sampled at an
instant. The current is low-pass filtered in analogue form before it reaches the converter, so
each stored value carries a weighted average of the recent current and not a reading of it at a
point. Writing the sampling period as $\Delta$ and taking the weighting to be uniform across it,
the quantity a likelihood must describe is the windowed average
\begin{equation}
  \bar{y}_t \;=\; \frac{1}{\Delta}\int_{t-\Delta}^{t} y(\tau)\,\mathrm{d}\tau \;+\; \eta_t,
  \label{eq:obs-avg}
\end{equation}
where $y(\tau)$ is the instantaneous population current and $\eta_t$ is the instrumental
noise accumulated over the window. This interval average is the physical observable. When
$\Delta$ is much shorter than the fastest relaxation of $\mathbf{Q}$ the average sits close
to the instantaneous current and the distinction is negligible; when $\Delta$ is comparable
to or longer than that relaxation the two diverge, and a likelihood built on instantaneous
sampling describes the wrong random variable.

Uniform weighting is a first approximation to what the electronics do. The analogue stage of a
patch-clamp amplifier is normally a Bessel low-pass, chosen because its group delay is nearly
flat and its step response barely overshoots, so the effective measurement kernel is that
filter's impulse response, a smooth shape whose tails reach into the neighbouring intervals.
The uniform window keeps the entire kernel inside one acquisition interval, which is what lets a
likelihood conditioned on the two ends of that interval account for the averaging without
reference to anything outside it. A Bessel kernel breaks that containment twice over: it
correlates successive recorded values through the signal, and it colours the instrumental noise
that Eq.~\ref{eq:obs-avg} takes to be white. We keep the uniform window throughout, in the
simulator as well as in every likelihood (Methods), so that data and likelihood describe the same
observable and the closures below are the approximation the diagnostic sees. The uniform window is
itself an approximation to the impulse response of the anti-aliasing filter of a real recording, and
the Discussion prices it; it does not enter the measurement here because the simulator and every
member share it.
% CORRECTED 2026-08-04. The clause read "the closures below are the ONLY approximation under test",
% which is a rule-4 universal and is contradicted by 05_discussion.tex, where the uniform window
% against a Bessel kernel is conceded as a second approximation adopted knowingly. What keeps that
% one out of the measurement is that it is shared, which is what the sentence now says. Carrying a general
kernel, with the coloured noise that comes with it, through the same construction is left to
later work.

\subsection*{Two Gaussian approximations}

Each member of the family makes two closures, on two different objects. To close a
distribution is to replace it by the Gaussian with the same mean and covariance,
discarding every higher moment. The word is borrowed from moment hierarchies, where that
substitution lets a finite set of equations stand for an infinite one, and it fits here
for the same reason: conditioning the occupancies on an observation takes them out of the
multinomial family, and the exact conditional distribution after $t$ samples has no
summary of fixed size, whereas $(\bm{\mu},\bm{\Sigma})$ is the same size at every step.

\begin{tcolorbox}[title={The two Gaussian approximations}, colback=gray!5, colframe=gray!55]
\textbf{Occupancy closure.} The joint distribution of the channel occupancies is exactly
multinomial when the channels are independent, which is also the maximum-entropy
distribution consistent with the mean occupancies. It is replaced by a multivariate
Gaussian with mean $\bm{\mu}$ and covariance $\bm{\Sigma}$, valid for large
$N_{\mathrm{ch}}$ by the central limit theorem and degrading as the channel count falls.

\smallskip
\textbf{Interval-signal closure.} The distribution of the interval-averaged conductance
over one window is a telegraph-like object with no closed form. It is replaced by a
Gaussian, valid when the window spans many transitions or when instrumental noise
dominates, and degrading as the window shrinks below the relaxation time.
\end{tcolorbox}

The higher moments these closures discard do not vanish. They reappear in the data as
temporal correlation that the likelihood does not predict, which is what the correlation
term of the information-distortion decomposition measures (Diagnostics). Three regimes
follow, and we name them here so the maps in the Results have vocabulary: the
\emph{multinomial} regime (few channels, where the occupancy closure fails), the
\emph{telegraphic} regime (windows far shorter than the relaxation time, where the
interval-signal closure fails), and the \emph{Gaussian} regime, where both closures hold
because the channel count is high enough for the occupancy Gaussian while the window is long
enough, or the instrumental noise large enough, for the interval-signal Gaussian. Any Gaussian likelihood of this
kind must fail somewhere. The aim of this paper is to show that the failures are the
predicted degradation of these two closures rather than defects of a particular algorithm.

The usage map of the Results is drawn on the plane of channel count against instrumental
noise at one acquisition interval, and its region names follow this vocabulary. % src: papers/_program/nomenclature.md (regions of the usage map)
Region~0, where the closure itself is misspecified, is the multinomial regime seen on that
plane, and its boundary carries a negative slope because instrumental noise makes the
emission more Gaussian and so relaxes the occupancy closure. Regions~1 to~4 (the unitary
current and channel count with a calibrated interval; the rates with a calibrated interval;
least squares at par; nothing estimable) all sit inside the Gaussian regime, and what
orders them is how much the data carry rather than whether the closures hold. The
telegraphic regime is a statement about $\Delta$ and so does not appear as a region of that
plane; it is reached by moving along the interval axis of Figure~\ref{fig:plane}.

\subsection*{The root question, and the ladder below it}

Before any of the choices above comes a prior question: does the method model the gating
fluctuations at all? Classical nonlinear least squares (\texttt{LSE}) answers no. It fits
the deterministic mean current and assigns one constant variance to every residual, so it
carries neither an occupancy distribution nor a conductance conditioned on anything, and
the temporal structure the two closures approximate is absent from it. \texttt{LSE} is
therefore not a member of the family. It is the bottom rung of the ladder of computational
cost this paper walks, and it is the anchor against which every member above it is
measured, because it is the method most readers are using.

Given a yes to the root question, the members differ along two axes. One axis is the
window: how much of the acquisition interval the conductance is conditioned on. The other
is recursion: whether the occupancy distribution is conditioned on the data as the
recording proceeds.

Along the recursion axis, a recursive (prequential) member factorizes the likelihood as
$\prod_t q_t(\bar{y}_t \mid \bar{y}_{1:t-1})$, so the occupancy distribution carried into
each window is the posterior left by the previous windows and the hidden state is
conditioned on the data. A non-recursive (open-loop) member factorizes as
$\prod_t m_t(\bar{y}_t)$, with each window's predictive propagated from the initial
condition and never conditioned on the observed trace.

Along the window axis the members form a ladder, indexed by how many of the interval's
endpoints the conductance is conditioned on (Table~\ref{tab:ladder}).

\begin{table}[htbp]
\begin{fullwidth}
\centering
\begin{tabularx}{\linewidth}{@{}ll>{\raggedright\arraybackslash}X@{}}
\toprule
Conditioned on & Members & Conductance model \\
\midrule
the fluctuations are not modelled & \texttt{LSE} & the deterministic mean current only \\
no endpoints & \texttt{NR}, \texttt{R} & instantaneous; the acquisition averaging is ignored \\
one endpoint (the start) & \texttt{MR}, \texttt{VR} & interval-mean conductance given the start state \\
endpoints not distinguishable & \texttt{INR} & interval-mean conductance; with no update the two endpoint rows coincide \\
two endpoints (the boundary) & \texttt{IR} & interval-mean conductance given both boundary states; interior marginalized \\
the full trajectory & (exact) & intractable; supplied by the stochastic simulation as ground truth \\
\bottomrule
\end{tabularx}
\end{fullwidth}
\caption{The window axis as an endpoint ladder, with least squares (\texttt{LSE}) shown at
the foot of the cost ladder although it lies off the lattice. Crossing the window axis with
the recursion axis gives the implemented grid: \texttt{NR} (non-recursive,
$\mathrm{av}=0$), \texttt{INR} (non-recursive, $\mathrm{av}=1$, total start-conditioned
interval variance), \texttt{R} (recursive, $\mathrm{av}=0$), \texttt{MR} (recursive,
$\mathrm{av}=1$, total start-conditioned interval variance), \texttt{VR} (recursive,
$\mathrm{av}=1$, residual interval variance), \texttt{IR} (recursive, $\mathrm{av}=2$),
where $\mathrm{av}$ counts the conditioned endpoints. \texttt{INR} occupies a row of its own
because without a recursive update the boundary state is unobservable: the two-endpoint
conductance can enter only through its start-state marginal, so $\mathrm{av}=1$ and
$\mathrm{av}=2$ describe the same non-recursive method. The body walks
\texttt{LSE}, \texttt{NR}, \texttt{R}, \texttt{IR}. Of the three reported in the Supplement,
\texttt{MR} and \texttt{VR} split the step from \texttt{R} to \texttt{IR}. MacroIR is
\texttt{IR} and MacroINR is \texttt{INR}.} % src: papers/_program/nomenclature.md (the lattice as implemented; the ladder for the Methods presentation). TODO 2026-07-31: the old clause "MNR is not separable from NR over the range measured" was REMOVED — it was measured on the defective build (missing N*ms); restate only after the 2026-07-31 INR re-run.
\label{tab:ladder}
\end{table}

All seven rungs are implemented and measured, and Figure~\ref{fig:time} carries them side by side under one set of conditions, so the lattice is read as a design rather than as a list.

Two things follow from the ladder. First, \texttt{IR} occupies the top rung below the
intractable exact likelihood, so its standing in the family is structural and set before
any measurement. Second, the lower rungs are established likelihoods rather than
constructed foils: the non-recursive instantaneous member is the independent-interval
likelihood of \citet{milescu2005maximum}, and the recursive instantaneous member is the
macroscopic recursive filter of \citet{moffatt2007estimation} and the generalized filter
of \citet{munch2022bayesian}. The lattice is one coordinate system that contains these
methods as special cases.

\subsection*{The boundary state}

A \emph{boundary state} is the pair $(i_0, i_t)$ of a channel's states at the two ends of
an acquisition interval. Chaining intervals, the end state of one interval is the start
state of the next, so there is a single state variable per junction: the filter conditions
on the states at the interval boundaries and marginalizes the trajectory between them
analytically. This is the same device as static condensation, or the spatial Markov
property, applied on the time axis. It is not a transition state in the mechanistic sense,
and the word \emph{transition} is not used for it here. Because the joint dynamics of the two
endpoints are generated by the Kronecker sum
$\mathbf{Q}\oplus\mathbf{Q}=\mathbf{Q}\otimes\mathbf{I}+\mathbf{I}\otimes\mathbf{Q}$, whose
propagator factorizes as $e^{(\mathbf{Q}\oplus\mathbf{Q})t}=e^{\mathbf{Q}t}\otimes
e^{\mathbf{Q}t}$, the moments over the $K^2$ boundary pairs collapse to bilinear forms in
the $K$-dimensional state space and no explicit $K^2$ object is formed (Methods).

The ladder needs two conductance objects, both built from the boundary-conditioned
single-channel moments $\overline{\Gamma}_{i_0\to i_t}$ (the mean interval-averaged
conductance of a channel that starts in $i_0$ and ends in $i_t$) and
$\overline{V}_{i_0\to i_t}$ (its variance). Both are moments of the integrated conductance of a
single channel over one window,
\begin{equation}
  A(\Delta)=\int_0^{\Delta}\gamma\bigl(X_s\bigr)\,\mathrm{d}s,
  \label{eq:integrated-conductance}
\end{equation}
taken conditional on the chain starting in $i_0$ and ending in $i_t$, and both are available in
closed form without diagonalizing $\mathbf{Q}$ and without forming any $K^2$ object. Write
$\bm{\Gamma}=\mathrm{diag}(\bm{\gamma})$ and build the two augmented generators
\begin{equation}
  \mathcal{Q}_2=\begin{pmatrix}\mathbf{Q} & \bm{\Gamma}\\ \bm{0} & \mathbf{Q}\end{pmatrix},
  \qquad
  \mathcal{Q}_3=\begin{pmatrix}\mathbf{Q} & \bm{\Gamma} & \bm{0}\\
                               \bm{0} & \mathbf{Q} & \bm{\Gamma}\\
                               \bm{0} & \bm{0} & \mathbf{Q}\end{pmatrix},
  \label{eq:augmented}
\end{equation}
of sizes $2K\times 2K$ and $3K\times 3K$. Their exponentials are block triangular with
$e^{\mathbf{Q}\Delta}$ on the diagonal, and the two blocks above it are the objects wanted:
$e^{\mathcal{Q}_2\Delta}$ has $\mathbf{F}_1(\Delta)$ in its $(1,2)$ block and
$e^{\mathcal{Q}_3\Delta}$ has $\mathbf{F}_2(\Delta)$ in its $(1,3)$ block, where
$\mathbf{F}_1$ and $\mathbf{F}_2$ are the first and second derivatives of the tilted semigroup
$e^{(\mathbf{Q}+\alpha\bm{\Gamma})\Delta}$ at $\alpha=0$, the second up to the factor
$\mathbf{W}=2\mathbf{F}_2$ that comes from the two triangular halves of the double integral.
This is the block construction of \citet{vanloan1978computing}. Dividing by the transition
probability conditions on the endpoints and dividing by $\Delta$ turns the integral into an
average, so
\begin{equation}
  \overline{\Gamma}_{i_0\to i_t}
    = \frac{[\mathbf{F}_1(\Delta)]_{i_0 i_t}}{\Delta\,P_{i_0\to i_t}(\Delta)},
  \qquad
  \overline{V}_{i_0\to i_t}
    = \frac{2[\mathbf{F}_2(\Delta)]_{i_0 i_t}}{\Delta^2\,P_{i_0\to i_t}(\Delta)}
      - \overline{\Gamma}^{\,2}_{i_0\to i_t}.
  \label{eq:gammabar-vbar}
\end{equation}
Both expressions are $0/0$ where a transition is unreachable, which is not a numerical edge case
here: at zero agonist $k_{\mathrm{on}}[A]$ vanishes identically, so $P_{C\to O}(\Delta)=0$ exactly
over the whole pre-pulse segment. Conditioning on a transit that occupies a vanishing fraction of
the window leaves the chain at its two endpoint states for essentially equal times, so the limit is
the endpoint average,
\begin{equation}
  \overline{\Gamma}_{i_0\to i_t}\to\tfrac12(\gamma_{i_0}+\gamma_{i_t}),
  \qquad
  \overline{V}_{i_0\to i_t}\to\Bigl(\tfrac{\gamma_{i_0}-\gamma_{i_t}}{2}\Bigr)^{2},
  \label{eq:endpoint-limit}
\end{equation}
which depends on the conductance vector alone and introduces no free constant.
% GAP CLOSED 2026-08-04, and it was the one that stopped a replication attempt in four of its six
% stages. The manuscript previously said only that these two objects "are computed from Q and the
% recording filter (Methods)", and neither Methods nor Appendix 1 gave a formula, a method or a
% citation, while these two symbols carry the whole av=1 against av=2 distinction, the MR/VR/IR
% algebra, the G and H matrices of Appendix 1 and the MR-IR identity.
% Sources transplanted, both verified against their own derivations:
%   theory/macroir/notes/Gmean_ij_gvarij/gmean_ij_gvar_by_blocks.tex (the two blocks, the reads,
%     W = 2F_2, the two normalisations)
%   theory/macroir/notes/Gmean_ij_gvarij/bayesian_prior_regularization_of_Qdt.md (the endpoint
%     average as the small-P limit, and why the smooth 1/sqrt(P^2+eps^2) replacement is worse:
%     it smuggles in a zero-conductance prior without flagging it)
% What stays out of the body and belongs in the Supplementary Information: the derivation of why
% F_1 and F_2 are those derivatives, and the argument for the endpoint average. Recipe here,
% derivation there.
The start-conditioned mean conductance is the $K$-vector
\begin{equation}
  (\bar{\gamma}_0)_{i_0} \;=\; \sum_{i_t} P_{i_0\to i_t}(\Delta)\,\overline{\Gamma}_{i_0\to i_t},
  \label{eq:gammabar}
\end{equation}
which conditions on the state at the interval's start alone and is the object used at
$\mathrm{av}=1$. The boundary-conditioned mean conductance is the same average retained on
both endpoints, so it stays a $K\times K$ object rather than collapsing to a vector, and is
the object used at $\mathrm{av}=2$.

\subsection*{The interval update}

Over a window of length $\Delta$, the occupancy mean and covariance propagate as
\begin{align}
  \bm{\mu}^{\mathrm{prop}} &= \mathbf{P}(\Delta)^\top \bm{\mu}_0, \label{eq:mu-prop}\\
  \bm{\Sigma}^{\mathrm{prop}} &= \mathbf{P}(\Delta)^\top\bigl(\bm{\Sigma}_0-\mathrm{diag}(\bm{\mu}_0)\bigr)\mathbf{P}(\Delta) + \mathrm{diag}(\bm{\mu}^{\mathrm{prop}}), \label{eq:sig-prop}
\end{align}
where $(\bm{\mu}_0, \bm{\Sigma}_0)$ is the occupancy Gaussian at the start of the window.
Every occupancy moment in this section is \emph{per channel}: with $\mathbf{N}_0$ the vector of
channel counts by state, $\bm{\mu}_0 = \mathbb{E}[\mathbf{N}_0]/N_{\mathrm{ch}}$ and
$\bm{\Sigma}_0 = \mathrm{Cov}(\mathbf{N}_0)/N_{\mathrm{ch}}$, so $\bm{\mu}_0$ is a probability
vector and $\bm{\Sigma}_0$ is a covariance divided by the channel count. The factors of
$N_{\mathrm{ch}}$ that appear below are what carry these onto the scale of the recorded current,
and they are the reason the rank-one correction of Eq.~\ref{eq:sig-post} carries one and the
correction to $\bm{\mu}$ does not.
% ADDED 2026-08-04. The convention was used consistently throughout this section and declared
% nowhere, which is worse than a contradiction because nothing in the text looks wrong. A reader
% implementing from Eq. mu-prop and Eq. sig-prop builds Sigma_0 as a total covariance and then
% every N_ch factor downstream is off by that factor. The source states it and flags it:
% theory/macroir/docs/Macro_IR/macroir_derivation.tex:73-95, "This convention is crucial: all
% covariance matrices in the spec are per-channel covariances, i.e. total covariances divided by
% N_ch." Found by a replication audit.
The predicted mean of the interval-averaged current is
\begin{equation}
  \bar{y}^{\mathrm{pred}} \;=\; N_{\mathrm{ch}}\,\bm{\mu}_0^\top \bar{\bm{\gamma}}_0,
  \label{eq:ypred}
\end{equation}
and its predicted variance decomposes into three independent contributions,
\begin{equation}
  \sigma^2_{\mathrm{pred}} \;=\; \epsilon^2_{\Delta}
  \;+\; N_{\mathrm{ch}}\,\widetilde{\bm{\gamma}^\top\bm{\Sigma}\bm{\gamma}}
  \;+\; N_{\mathrm{ch}}\,\bm{\mu}_0^\top\bm{\sigma}^2_{\bar{\gamma}_0}.
  \label{eq:pred-var}
\end{equation}
The first term $\epsilon^2_{\Delta}$ is the instrumental noise over the window, which
carries the acquisition averaging and decreases with $\Delta$; its parametric form is given
in Methods. The second term is the population gating variance, the fluctuation of the window
current transmitted through the occupancy covariance. The tilde on it is an operator rather
than a decoration, and the symbols under it are not the $K$-vector $\bar{\bm{\gamma}}_0$ and
the $K\times K$ matrix $\bm{\Sigma}$ used elsewhere in this section: it denotes a bilinear
form taken over boundary \emph{pairs}, contracted back to a scalar. Writing
$\bm{\Sigma}^{\mathrm{bnd}}_{0\to\Delta}$ for the per-channel covariance of the boundary-pair
occupancies,
\begin{align}
  \widetilde{\bm{\gamma}^\top\bm{\Sigma}\bm{\gamma}}
  &\;=\; \sum_{i_0,i_t}\ \sum_{j_0,j_t}
    \overline{\Gamma}_{i_0\to i_t}\,
    \bigl(\bm{\Sigma}^{\mathrm{bnd}}_{0\to\Delta}\bigr)_{(i_0,i_t),(j_0,j_t)}\,
    \overline{\Gamma}_{j_0\to j_t} \nonumber\\[2pt]
  &\;=\; \bar{\bm{\gamma}}_0^\top\bigl(\bm{\Sigma}_0-\mathrm{diag}(\bm{\mu}_0)\bigr)\bar{\bm{\gamma}}_0
    \;+\; \bm{\mu}_0^\top\bigl(\mathbf{H}\mathbf{1}\bigr),
  \label{eq:tilde}
\end{align}
where $H_{i_0 i_t} = \overline{\Gamma}^2_{i_0\to i_t}\,P_{i_0\to i_t}(\Delta)$ and $\mathbf{1}$
is the vector of ones. The first line is the definition and the second is what is evaluated:
only $K$-vectors and $K\times K$ matrices appear, so the $K^2\times K^2$ boundary covariance
is never formed. The distinction matters for reading Eq.~\ref{eq:pred-var}, because taking the
tilde to mean $\bar{\bm{\gamma}}_0^\top\bm{\Sigma}_0\bar{\bm{\gamma}}_0$ keeps the first
summand and drops the second, which is the single-channel spread of the interval mean across
boundary pairs and does not vanish when the occupancy covariance does. The general operator,
of which this is the bilinear case, is given in Appendix~1. The third term of
Eq.~\ref{eq:pred-var} is the within-window per-channel variance, with
$\bm{\sigma}^2_{\bar{\gamma}_0}$ the start-conditioned interval variance built from
$\overline{V}_{i_0\to i_t}$.

Writing $\mathbf{g}=\mathrm{Cov}(\text{end-of-window occupancy}, \bar{y})/N_{\mathrm{ch}}$
for the per-channel cross-covariance between the interval current and the occupancy at the
end of the window, and $\delta = \bar{y}^{\mathrm{obs}} - \bar{y}^{\mathrm{pred}}$ for the
innovation, the Gaussian update of the occupancy is the scalar Kalman correction
\begin{align}
  \bm{\mu}^{\mathrm{post}} &= \bm{\mu}^{\mathrm{prop}} + \frac{\mathbf{g}}{\sigma^2_{\mathrm{pred}}}\,\delta, \label{eq:mu-post}\\
  \bm{\Sigma}^{\mathrm{post}} &= \bm{\Sigma}^{\mathrm{prop}} - \frac{N_{\mathrm{ch}}\,\mathbf{g}\mathbf{g}^\top}{\sigma^2_{\mathrm{pred}}}. \label{eq:sig-post}
\end{align}
The factor $N_{\mathrm{ch}}$ enters the covariance downdate but not the mean correction
because $\bm{\Sigma}$ is carried on the population scale, the same $N_{\mathrm{ch}}$ that
multiplies the gating variance in Eq.~\ref{eq:pred-var}, whereas $\bm{\mu}$ and the gain
$\mathbf{g}$ are per-channel objects, so the rank-one correction to $\bm{\Sigma}$ must be
rescaled by $N_{\mathrm{ch}}$ to sit on its scale while the correction to $\bm{\mu}$ needs
no such factor.
In a recursive member the posterior pair $(\bm{\mu}^{\mathrm{post}}, \bm{\Sigma}^{\mathrm{post}})$
becomes the prior for the next window; in a non-recursive member it is not carried forward,
and each window's predictive is propagated from the initial condition alone. Reusing the
posterior as the next prior is itself a closure: it is exact only while the end marginal is
a sufficient summary of the boundary posterior, which the Gaussian projection does not
guarantee. For this reason the recursive members' Gaussian information is a model-based
quantity and is not claimed to be the exact information of the process; how far it departs
from the truth is what the Diagnostics measure.

Two algebraic differences separate the members that sit at $\mathrm{av}=1$ and
$\mathrm{av}=2$. The first is the form of the third term of Eq.~\ref{eq:pred-var}.
\texttt{MR} carries the total per-start-state interval variance, the spread of the
interval-averaged conductance of a channel known only to have started in $i_0$, while
\texttt{VR} and \texttt{IR} carry the residual variance left once the end state is also
conditioned on, $\sum_{i_t} P_{i_0\to i_t}(\Delta)\,\overline{V}_{i_0\to i_t}$. By the law
of total variance the two forms differ by the spread of the interval mean across end
states,
\begin{equation}
  \underbrace{\sum_{i_t} P_{i_0\to i_t}(\Delta)\bigl[\overline{V}_{i_0\to i_t}
    + \bigl(\overline{\Gamma}_{i_0\to i_t}-(\bar{\gamma}_0)_{i_0}\bigr)^2\bigr]}_{\text{total, }\texttt{MR}}
  \;=\; \underbrace{\sum_{i_t} P_{i_0\to i_t}(\Delta)\,\overline{V}_{i_0\to i_t}}_{\text{residual, }\texttt{VR},\,\texttt{IR}}
  \;+\; \mathrm{Var}_{i_t}\bigl[\overline{\Gamma}_{i_0\to i_t}\mid i_0\bigr],
  \label{eq:mr-vr}
\end{equation}
so \texttt{MR} charges to the observation variance a spread that the end state, once
conditioned on, would resolve. The second difference is the boundary cross-covariance term
\begin{equation}
  N_{\mathrm{ch}}\,\widetilde{\bm{\gamma}^\top\bm{\Sigma}\bm{\gamma}}
  \label{eq:mr-ir}
\end{equation}
that enters the gain $\mathbf{g}$ of Eqs.~\ref{eq:mu-post} and~\ref{eq:sig-post}, which
\texttt{IR} retains and both $\mathrm{av}=1$ members drop. Stepping \texttt{MR} to
\texttt{VR} changes the variance form and nothing else, and stepping \texttt{VR} to
\texttt{IR} restores the boundary term in the gain and changes nothing else, % src: papers/_program/nomenclature.md (the two concrete differences that drive the ladder)
so the two halves of the \texttt{R} to \texttt{IR} step are separated and can be measured
one at a time. Their consequences for the predicted variance across the regime grid are
characterized in the Results.

\subsection*{The top of the ladder}

The exact likelihood conditions on the full trajectory inside each window and is
intractable, which is why the two closures above are needed at all. The same stochastic
model that makes the exact likelihood intractable can be simulated forward exactly, one
trajectory at a time, so the simulation realizes the object at the top of the ladder and
serves as the ground truth against which every rung below it is judged.
